\section{Implications}

Parts~\ref{pt:data-driven-station-planning} and \ref{pt:predictive_operations} have discussed in detail the methodological choices, results, and limitations of this thesis's contributions. This Section now places those findings within the broader context established in Parts~\ref{pt:introduction} and~\ref{pt:foundations}, considering how this work addresses—or does not address—the overarching questions and challenges.

\subsection{From context to implications}

Part~\ref{pt:introduction} presented a reality of urban environments under pressure. Cities concentrate a large share of energy use and emissions, and urban mobility--long dominated by car-centric planning--still feeds \acp{ghg}, congestion, and inequity, despite growing investment in \ac{pt}, cycling, and shared mobility. Furthermore, this pressure is expected to intensify in the coming years~\cite{UN_SDG11_2015}, with clear implications for the environment and human health. While recent decades have brought advances in public policy, shifts in public attitudes toward sustainability, and considerable developments in technology, the fundamental urban challenge remains. Car dependency, emissions, and inequitable access to transport persist as unresolved issues, showing that further efforts are needed to solve the overarching problem.

\acp{bss}, as a low-carbon, flexible, and active mobility option, stand out as one of the most visible responses to this challenge, and their popularity and visibility have grown in recent years. Despite this success, there is still much to do. Ongoing work is needed so that more people not only use but also see these more eco-friendly mobility options as a viable alternative to car-based travel, and at the same time, to further reduce the negative impact of these systems on the environment and human health. Only through sustained efforts to make such systems more efficient, equitable, and compelling can their positive impact on urban sustainability grow.

This thesis contributes to these needs by developing computational data-driven methods, using the kind of operational and open urban data that many cities and \ac{bss} operators already collect. It can therefore be seen as a step along the journey that DeMaio~\cite{DeMaio2009} and Shaheen et al.~\cite{Shaheen2010} envisioned for the fourth generation of \acp{bss}: systems that are more efficient, sustainable, and user-friendly through improved rebalancing, the integration of \acp{e-b}, and data-driven management.


\subsection{Broadening the user base}

To address the earlier challenges, sustainable mobility options like \acp{bss} must attract a broader range of users. This would help reduce \acp{ghg} emissions and car dependency and increase mobility opportunities, especially for those without private vehicles. In this line, Section~\ref{sec:who_uses_bss} highlights two main points. First, \acp{bss} users are typically younger, white, male, well-educated, and have higher incomes than the general population~\cite{Murphy2015, BachandMarleau2012, Goodman2014, Shaheen2014, Fishman2013, Fishman2014, Bauman2017, Ogilvie2012, Steinbach2011, Willberg2021, Waldner2025}. Second, \ac{bss} usage depends heavily on station locations, and systems are usually deployed in central, higher-income areas~\cite{Ogilvie2012, Goodman2014, Willberg2021}. These factors have led some to describe bike-sharing as one of the most inequitable forms of sustainable transportation infrastructure~\cite{Hoffmann2016}. Therefore, there is potential to broaden the user base.

Because the spatial distribution of stations and system design influence participation, expanding access requires tools that help planners test planning scenarios beyond aggregate ridership alone. The station-planning framework developed in this thesis is designed to address exactly these challenges. It enables planners to optimize for a variety of objectives and to balance trade-offs between competing goals, all within a single tool. For instance, as discussed in Section~\ref{sec:who_uses_bss}, women tend to be more risk-averse than men~\cite{Garrard2008, Garrard2012}, and in London, they were more likely to make bike-sharing trips that started or ended in parks, suggesting a preference for routes that avoid motorized traffic and favor quieter, traffic-free environments~\cite{Goodman2014, Johnson2010}. With this framework, planners could prioritize station placement in areas with dedicated cycling infrastructure or car-free areas, helping make the system more attractive to women.

Moreover, socioeconomic disparities play a major role in shaping access to bike-sharing, as most \acp{bss} have historically been concentrated in central, higher-income neighborhoods, leaving outlying and lower-income areas under-served~\cite{Ogilvie2012, Goodman2014, Willberg2021}. However, studies have shown that residents of these underserved neighborhoods are willing to use bike-sharing when it is made available~\cite{Yu2024}. The framework developed in this thesis directly supports equity-oriented planning goals by enabling planners to evaluate and prioritize more inclusive scenarios. By incorporating data on car ownership or vehicle availability, the framework also helps planners design systems that better serve those with fewer transportation options. Evidence indicates that individuals without car access are more likely to use \acp{bss}~\cite{Mohiuddin2024, Yu2024}, so expanding the system into neighborhoods with lower rates of car ownership can help address unmet mobility needs in those communities. Additionally, there is evidence that placing stations in targeted neighborhoods can have the extra benefit of modestly decreasing household vehicle ownership~\cite{Basu2021}. Finally, topography plays a crucial role in hilly cities~\cite{Midgley2011, MateoBabiano2016, Morency2015}. This framework includes a topography module, enabling planners to allocate more stations to steeper areas and thus distribute accessibility more fairly across the city. This consideration is especially important when \acp{e-b} are not available to mitigate the challenges posed by slopes.

On the other hand, age is not a factor that can be ignored. Older adults remain underrepresented in many \acp{bss}~\cite{Yu2024, Willberg2021}, and good station planning cannot address this issue by itself. The introduction of \acp{e-b} in the 4th generation of \acp{bss}, however, has made available to this group a mobility option that is more accessible and attractive, reducing physical effort on long or hilly trips. Recent evidence from the Netherlands shows that \acp{e-b} have led to higher usage among older women, even narrowing the gender gap among older cohorts~\cite{Huang2024}. This thesis contributes to making \acp{e-b} a reliable option where they are needed: maintenance and battery predictions help operators maintain bike availability and rebalance the fleet. 

Trust and perceived reliability also influence who uses \acp{bss}. Operational problems can quickly undermine public confidence~\cite{Yle2015_HelsinkiBikeShare}. When stations are frequently empty or full, or when bikes are often out of service or poorly maintained, potential users are discouraged. Section~\ref{sec:operational_indicators} summarizes some of the indicators operators already use--\ac{tdb}, trips per 1,000 residents, station non-operational time, dock occupancy, maintenance performance, and, with electrification, battery state of charge~\cite{Waldner2025, Yanocha2018}--and stresses that maintenance is crucial for long-term sustainability~\cite{Rossow2018, Santander2022}. This thesis adds two interpretable indicators that can be considered when monitoring a system: \textbf{inter-event time} at a station (the time between consecutive trips) and a \textbf{distance-based index} (combining trip distances and origin frequencies). Both moderately correlate with station-level average battery and help identify where depletion is likely to be fastest, so operators can target recharging, redistribution, or replacement without requiring new sensor infrastructure. In parallel, the failure-risk estimates from the maintenance contribution support preventive rather than purely reactive repairs, improving the likelihood that bikes are in good working order when users need them. Better availability and a better state of the fleet can thus contribute to the kind of reliable service that Part~\ref{pt:foundations} associates with greater uptake.

Taken together, these contributions do not by themselves broaden the user base, but they give planners and operators concrete levers: more equitable station placement, more reliable \ac{e-b} availability, and better fleet condition, which can make bike-sharing a more credible and attractive option for a wider range of users.


\subsection{System design and management}

Part~\ref{pt:foundations} has shown how \acp{bss} are configured and run: where stations go, which bikes they offer, how they are maintained and rebalanced, and who tends to use them. These choices are not neutral. They shape who has access, how much the system is used, and what it costs to operate. This Section considers what the thesis contributes to the way systems are designed and managed in practice.

Designing a \ac{bss} is often seen as a matter of maximizing demand or geographic coverage. Yet, experience demonstrates that these are not the only priorities~\cite{Yanocha2018}. In practice, planners frequently need different tools for different objectives. This thesis addresses this gap by introducing a single planning framework that can accommodate multiple, potentially conflicting objectives, allowing systems to be designed and optimized according to diverse priorities~\cite{GrauEscolano2025}. 

When the objective is to maximize demand or coverage, Section~\ref{sec:factors_influencing_bss_usage} outlines a range of relevant inputs: it summarizes determinants of \ac{bss} usage from broader contextual and system-level design choices to specific station-level factors, such as land-use mix, density, \ac{pt} proximity, cycling infrastructure, and station density. Importantly, the framework allows planners not only to incorporate these determinants explicitly, but also to prioritize among them, enabling more precise scenario testing and tailored planning.

Placing stations in less favored areas in terms of demand can also be desirable, for two reasons. First, \acp{bss} do not only follow demand; they can generate it. Chapter~\ref{ch:impacts_of_bss} shows that \acp{bss} can stimulate local commerce, trips and spending, and property-market effects~\cite{Schoner2012, Buehler2014, ElGeneidy2016}, so positioning stations where new activity is desired can help promote those areas and attract more people. Second, where stations go shapes who benefits, and if deployment follows only existing demand, spatial inequalities can be reinforced. Extending coverage to less favored areas is therefore a matter of equity as well as of local development. The framework supports this aim by allowing equity-oriented scenarios (e.g.\ prioritizing low-income population or areas with low vehicle availability within the service area) and by including a topography module so that accessibility can be distributed more evenly in hilly cities, where slope would otherwise exclude many users unless mitigated by \acp{e-b}. Planners can thus use the same tool to pursue demand and coverage, to stimulate less favored areas, or to trade off between these or other goals.

Another system design choice is which bikes to include. For commuting-oriented systems, \acp{m-b} and \acp{e-b} are the usual options, and in other contexts, cargo bikes or other types may be more appropriate. The decision should be informed by the potential user base and trip purposes. This thesis adds that the choice has consequences later in the system life: \acp{e-b} and \acp{m-b} exhibit distinct mobility patterns and distinct component-level failure dynamics~\cite{Grau-Escolano2024}. So the mix of bike types chosen at design time will shape rebalancing needs, maintenance scheduling, and spare-parts planning throughout operation.

The thesis also contributes two use cases that aim to facilitate operation. Predicting \ac{e-b} locations and battery levels supports rebalancing and recharging: short-horizon forecasts help operators target where bikes and charge are needed, and simple indicators such as inter-event time and a distance-based index help identify stations where depletion is probably fastest~\cite{Bassolas2025}. Predicting component-level failure risk supports maintenance: the survival models provide interpretable, fleet-wide estimates that can inform workshop scheduling and spare-parts planning, so that repairs can be preventive rather than purely reactive~\cite{Grau-Escolano2024}. Together, these tools can reduce the operational burden that follows from a given design and make it easier to sustain service quality where it matters most.

Ultimately, cost is the primary constraint in the design and operation of a \ac{bss}. Both the system's configuration and its daily running must be balanced against financial viability. Ideally everyone would have a station close by, but that is rarely feasible. An excessive number of stations greatly increases operational expenses—especially for rebalancing—while too few, or stations placed too far apart, make the system inconvenient and inaccessible for users. Despite most docked \acp{bss} being publicly owned or managed through public-private partnerships~\cite{Li2021, Cao2022}, and many depending on subsidies~\cite{Shaheen2010}, there is still an expectation that they should be broadly viable. While the tools introduced here may be valuable for addressing specific aspects, designing a \ac{bss} at the system-wide level demands a comprehensive assessment of costs—a perspective that lies beyond the intended scope of this thesis.


\subsection{Environmental impacts}

As discussed in Section~\ref{sec:environmental_impacts}, manufacturing, rebalancing, maintenance, and battery charging all contribute to the life-cycle footprint~\cite{Luo2019, Calan2024, Zhu2023}. In docked systems, manufacturing is often the largest single contributor, accounting for around 50\% of total life-cycle emissions, while rebalancing alone can account for 36\% of life-cycle \acp{ghg}. Maintenance, when included, often represents around 10--15\% and is sensitive to assumptions about component lifetimes~\cite{Luo2019, Calan2024, Zhu2023, Sun2022}. Moreover, environmental benefits from avoided transport emissions materialize primarily when car substitution is substantial~\cite{Luo2019}.

The \ac{pm} contribution in Chapter~\ref{ch:failure_forecasting_bicing} can both improve how the environmental impacts of \acp{bss} are computed and offer a way to reduce them~\cite{Grau-Escolano2024}. It first provides a description of maintenance patterns for three key components by bike type based on real operational data. Such information can improve \ac{lca} assumptions about replacement intervals and component life where it is missing. Since \acp{lca} often rely on generic or operator-average values, sharing this evidence across systems could build a shared knowledge base that, even if not perfectly transferable across systems, could still be valuable. The bike component failure predictions, in turn, enable predictive rather than corrective maintenance. This can helps extend component lifespans and reduce unnecessary replacements, which lowers demand for new parts. Given manufacturing's share of life-cycle emissions, extending component lifetimes can meaningfully reduce the overall environmental footprint of \acp{bss}.

Second, the battery forecasting presented in Chapter~\ref{ch:mobility_and_battery_modeling_bicing} allows operators to more efficiently plan recharging and rebalancing by predicting where bikes and charged batteries are likely to be needed~\cite{Bassolas2025}. More accurate system-level forecasts can, in theory, reduce the total distance traveled by operational vehicles for rebalancing, thereby lowering the use-phase emissions that \acp{lca} attribute to this activity--a particularly relevant point since rebalancing accounts for over one third of total \acp{ghg} emissions. However, electrification adds complexity to rebalancing: operators must consider not only bike availability but also battery levels, which can increase the frequency of interventions. While the battery forecasting tool supports more efficient management of these needs and may improve targeting, the overall effect on operational travel is uncertain, as electrified systems may still demand more rebalancing than purely mechanical ones.

Lastly, the station-planning framework in Part~\ref{pt:data-driven-station-planning} supports designs that favor scenarios such as maximizing demand coverage and first--last mile connectivity~\cite{GrauEscolano2025}. Such layouts can encourage system use and promote modal substitution away from the car, and therefore would reduce \acp{ghg}. The foundations show, however, that the magnitude of avoided emissions depends critically on realized substitution patterns and that benefits are often overestimated when substitution is assumed rather than measured~\cite{Kou2020, Saltykova2022, Liu2021}. Therefore, the impacts of any given design would require dedicated assessment.


\subsection{Health and safety impacts}

Section~\ref{sec:health_and_safety_impacts} shows that net health impacts of \acp{bss} are driven mainly by gains in physical activity, while traffic injury and air pollution play smaller roles in most settings~\cite{Rojas-Rueda2011, Woodcock2014, Otero2018, Teixeira2021}. Who benefits depends on who has access and \ac{hia} evidence suggests that prioritizing higher-poverty or less-served areas can yield larger aggregate benefits than deployment that follows existing demand~\cite{Babagoli2019, Clockston2021}. The size of those gains, however, depends on which modes trips replace~\cite{Fishman2014, Murphy2015, Bauman2017}: when new users mainly substitute walking or \ac{pt} rather than car travel, population-level physical activity gains can be limited. On the safety side, injury risks vary across cities and depend on local road-safety environments and infrastructure~\cite{Rojas-Rueda2011, Woodcock2014, Otero2018, Clockston2021}.

The link between this thesis and health and safety is mainly indirect: the tools developed here can help bring more people into the system. The station planning framework supports equity-oriented and topography-aware station distributions that extend access to underserved areas, in line with the conditions under which \ac{hia} suggest benefits can be larger, and it allows prioritizing locations with cycling infrastructure or lower traffic exposure~\cite{Goodman2014, Garrard2012}, which could favor safer environments. Moreover, the battery and maintenance contributions support more reliable \ac{e-b} availability and fleet condition. Although \acp{bss} have been shown to be relatively safe~\cite{Woodcock2014, Martin2016, Fishman2016_road_safety, Jacobsen2015, Teixeira2021}, \ac{pm} can further enhance user safety by lowering the chances of component failures during rides. Taken together, these advances enable the system to reach and retain a wider range of users, and thus, provide greater opportunities for the main health benefit of bike-sharing, increased physical activity, to be realized. How much of that potential is realized, and how the balance between health gains and risks from injury or pollution plays out, would depend on local context and on dedicated assessment of any given design~\cite{Teixeira2021, Chen2023}.


\subsection{Computational methods}

Computational methods have been central to the design and operation of \acp{bss} for over a decade (Chapter~\ref{ch:computational_methods_in_bss}). As the volume and complexity of generated data have grown, they have evolved from supplementary tools to core elements of how these shared mobility systems are run~\cite{FaghihImani2017, CortezOrdonez2024}. The data landscape available for such work has expanded considerably~\cite{Talavera_Garcia2021, Yanocha2018, Waldner2025}, and the literature has kept pace, with sustained progress in fields such as demand forecasting, rebalancing, and infrastructure planning. Nevertheless, there remains room for improvement in multiple areas. This thesis has focused on using data that remain seldom available or exploited in research—such as maintenance records and battery state—and on applying computational methods both to problems that are rarely addressed in the literature and to advancing the state of the art where methods already exist.

A central theme across the thesis is the importance of high-resolution and high-quality data for meaningful analysis and prediction in \acp{bss}. Across planning and operations, the results show that spatiotemporally detailed, contextual and operational data can reveal pronounced heterogeneities within the system. Different parts of the city, stations, and bicycles exhibit distinct patterns of accessibility, usage, energy depletion, and component wear that would be largely obscured under aggregated representations. Moreover, beyond supporting descriptive analyses, higher-resolution data can also lead to improved predictions by reducing the need for simplifying assumptions.

Beyond data, network effects play a key role in shaping how \acp{bss} function. The findings throughout this thesis show that system performance depends not just on individual station or trip features, but on how stations are positioned within the wider network and how they relate to one another. Because the network defines both station-to-station proximity and system accessibility, it shapes how users move between stations. These same network-driven flow patterns concentrate trips on certain station pairs and corridors, which in turn can lead to faster battery depletion and greater component wear.

The type of network used to represent the system is therefore important for the realism of the analysis and modeling. In this thesis, a representation grounded in the street layout proved effective: a directed graph in which distances are measured along the actual street network rather than in a straight line, and in which one-way streets and turn restrictions are respected, so that travel times and accessibility can differ by direction. Incorporating topology also proved important--a factor that is often overlooked in network-based analyses even though it shapes who can access stations, which routes users take, and how much energy trips consume. The combination of a directed, street-based network with topography was found to increase the realism of both the analysis and the modellings. 

The same concern for realism and practical use extends to the choice of models: the thesis favors methods that decision-makers can understand and trust, as those who use these tools need to see how inputs lead to outputs to be confident in their use. Chapter~\ref{ch:computational_methods_in_bss} shows a strong trend toward \ac{dl} and other high-capacity models such as \acp{gnn}, especially for demand forecasting. Such models can achieve high accuracy but typically need substantial data and compute, and their internal logic is hard to explain. When outcomes must be justified, simpler models are often preferable. Post-hoc techniques exist to explain \ac{dl} outputs, but where possible this thesis uses methods that are interpretable by design, so that users can explain and act on the results without extra tooling. That preference is reflected in each of the three contributions, though the way interpretability is achieved differs. In station planning, it comes from the chosen criteria and weights: planners see how their priorities map to layouts and can adjust them. In battery and location forecasting, it comes from transparent building blocks (Markov mobility, physics-based battery), so that inputs like trip flows, rest times, distance, and slope directly determine where charge is needed. In \ac{pm}, several options were compared, but the \ac{dl} variant \ac{deepsurv} performed best. Rather than give up that accuracy, the thesis applies post-hoc interpretability techniques (\ac{shap} and partial dependence analysis) so that the main drivers of failure risk remain visible.

These priorities, together with the structure and data of each problem, led to specific modeling choices. In station planning, finding high-quality station configurations is a combinatorial problem for which exact methods become prohibitive as the number of candidates and stations grows--as in the rebalancing and infrastructure literature (Chapter~\ref{ch:computational_methods_in_bss}). The thesis therefore relies on a metaheuristic, which scales to realistic instance sizes but does not guarantee global optimality, a trade-off typical of metaheuristic search. Among the options commonly used in \ac{bss} and facility-location problems a \ac{ga} was chosen and adapted to the structure of the problem. Other algorithmic choices would have been possible--for instance multi-objective algorithms such as \ac{nsgaii} or \ac{spea2}, which approximate Pareto fronts directly. Such methods would yield a set of non-dominated solutions in a single run, giving planners a direct view of the trade-offs between competing objectives without having to fix weights in advance or run multiple scenarios, and therefore, potentially simplifying the decision-making process. On the other hand, a novelty of the framework is to make spatial station layout an explicit object of the optimization. Where earlier \ac{bss} planning models typically treat the arrangement of stations as an incidental outcome of the optimization process, this one incorporates proximity and accessibility metrics so that planners can deliberately balance clustered versus dispersed configurations while still respecting the scenario-specific criteria and weights.

For the battery predictions, the thesis used trip records and system-level battery snapshots. \acp{gnn} and other \ac{dl} approaches could in principle be applied to predict station-level availability (as in the demand-forecasting literature) or battery status. However, the availability of battery data in terms of resolution and coverage was limited (e.g.\ 30-minute snapshots over only a few weeks, or 8-hour snapshots over a longer period). Therefore, data-intensive \ac{dl} models would be poorly suited to these constraints. Moreover, other options such as classical time series or tree-based \ac{ml} on tabular features would face the same problem. Therefore, this thesis opts instead for the hybrid (Markov mobility, physics-based battery, and simple indicators) noted above, which is data-efficient and interpretable given the data at hand. 

For the \ac{pm} problem the data includes when parts were replaced, but for many bikes the component was still in service at the end of the observation period. Using only observed replacements would therefore give a partial, pessimistic picture. For this kind of data, regression, classification, time series, and standard \ac{ml} or \ac{dl} approaches are ill-suited, as they typically require a single observed outcome per component, which is missing for units still in service. Therefore, survival models are the natural choice. Several \ac{sa} models were compared, and \ac{deepsurv} performed best, consistent with the broader trend toward \ac{dl}.

Beyond these methodological implications, the scope and transferability of the results deserve mention. The implications discussed above are drawn from analyses grounded in Barcelona and its \ac{bss}, \textit{Bicing}. The methods themselves are nevertheless transferable in principle to other urban contexts, as they rely on types of data that many \acp{bss} and cities already collect or can obtain.